% ID: FIS-CIR-OHM-002
% Titolo: Lampadina alimentata da un generatore
% Disciplina: Fisica
% Area: Circuiti elettrici
% Argomento: Legge di Ohm
% Sottoargomento: Applicazione numerica in un circuito semplice
% Classe: Prima istituto professionale
% Difficolta: 2
% Tipo: problema_numerico
% Risultato: soluzione da aggiungere
% Tempo_stimato: 8 min
% Competenze: applicazione della legge di Ohm, calcolo algebrico, interpretazione fisica
% Prerequisiti: tensione, corrente elettrica, resistenza, legge di Ohm
% Tag: fisica, circuiti_elettrici, legge_ohm, lampadina, generatore
% Fonte: verifica_fisica_circuiti_anonimizzata
% Autore: Riccardo Carli
% Licenza: CC-BY-SA-4.0

\begin{esercizio}
In un circuito elettrico elementare una lampadina ha una resistenza di
\(8\,\Omega\). Una corretta illuminazione necessita di una corrente elettrica
di \(0{,}5\,\mathrm{A}\).

Quale tensione deve sviluppare il generatore? Se la resistenza della lampadina
fosse la meta, quanta corrente passerebbe nel circuito mantenendo il valore di
tensione trovato in precedenza?
\end{esercizio}

\begin{soluzione}
% Soluzione da aggiungere.
\end{soluzione}

